\begin{helper}[A disjoint Kempe pair leaves the other incidence graph unchanged]
Let $\varphi$ be a partial edge coloring, and let
$\alpha,\beta,\gamma,\delta$ be colors such that the two finite color pairs
are disjoint:
\[
\{\alpha,\beta\}\cap\{\gamma,\delta\}=\varnothing.
\]
For any node $z$, swapping $\gamma$ and $\delta$ on the full
$(\gamma,\delta)$-component through $z$ leaves the entire
$(\alpha,\beta)$-incidence graph unchanged.
\end{helper}
\begin{proof}
Write $\varphi'$ for the coloring obtained from
\noderef{TwoColorComponentSwap}.  Consider an individual incidence between
a vertex $v$ and an exact edge identity $e$.  By
\noderef{TwoColorIncidenceGraph}, it is an edge of the
$(\alpha,\beta)$-incidence graph precisely when $v$ is incident with
$e$ and the color of $e$ is $\alpha$ or $\beta$.

If the edge-node of $e$ is outside the switched component, then
$\varphi'(e)=\varphi(e)$.  If it is inside, the swap either exchanges
$\gamma$ and $\delta$ or fixes the value of $e$.  Since neither
$\gamma$ nor $\delta$ equals $\alpha$ or $\beta$, the predicate
``the color is $\alpha$ or $\beta$'' has the same truth value before and
after the swap.  The incidence relation itself is unchanged.  Thus every
vertex--edge adjacency is identical in the two incidence graphs.  Both
graphs have no vertex--vertex or edge--edge adjacencies by definition, so
the graphs are equal.
\end{proof}
